%============================================================
% CH1 · P4 · PROBLEM 1 — TWO BUTTONS (Codeforces 520B)
%============================================================
\subsection{Problem 1.4.1: Two Buttons}
\problemlabel{1.4.1}{ssec:ch1-two-buttons}

\problemcard{1.4.1}{Two Buttons}{Codeforces 520B}{https://codeforces.com/problemset/problem/520/B}{\exmedium\ \cf{1400}}{Numbers as vertices, button presses as edges; BFS in a proven window; reverse greedy}

\subsubsection*{1. Problem statement}

A device shows a positive integer on its display. The \textbf{red} button multiplies the number by $2$, the \textbf{blue} button subtracts $1$. If the number ever stops being positive, the device breaks. The display can show arbitrarily large numbers. Initially the display shows $n$. What is the \textbf{minimum number of presses} needed to make it show $m$?

\textbf{Constraints:} $1 \le n, m \le 10^4$, \ $n \ne m$. \quad \textbf{Input:} one line with $n$ and $m$. \quad \textbf{Output:} the minimum number of presses.

\begin{table}[H]
\centering\small
\begin{tabular}{@{}cccl@{}}
\toprule
$n$ & $m$ & \textbf{Output} & \textbf{One optimal sequence} \\
\midrule
4 & 6 & 2 & blue, red: \ $4 \to 3 \to 6$ \\
10 & 1 & 9 & blue nine times: \ $10 \to 9 \to \dots \to 1$ \\
\bottomrule
\end{tabular}
\end{table}

\begin{diagrambox}
\begin{tikzpicture}
    % the device
    \draw[rounded corners=4pt, draw=deepblue, fill=theorybg, line width=1pt] (0,0) rectangle (4.2,2.4);
    \draw[rounded corners=2pt, fill=darkgray!85] (0.5,1.3) rectangle (3.7,2.1);
    \node[font=\Large\ttfamily\bfseries, text=successgreen!60] at (2.1,1.7) {4};
    \node[circle, draw=recursionred!80!black, fill=recursionred, minimum size=7mm, text=white, font=\scriptsize\bfseries] at (1.2,0.6) {$\times2$};
    \node[circle, draw=primaryblue!80!black, fill=primaryblue, minimum size=7mm, text=white, font=\scriptsize\bfseries] at (3.0,0.6) {$-1$};
    \node[font=\scriptsize, text=darkgray] at (1.2,-0.25) {red};
    \node[font=\scriptsize, text=darkgray] at (3.0,-0.25) {blue};
    % sample 1 as a path
    \begin{scope}[xshift=6cm, yshift=1.2cm]
        \node[gvsrc] (a) at (0,0) {4};
        \node[gv] (b) at (2,0) {3};
        \node[gvdst] (c) at (4,0) {6};
        \draw[gpath, -{Stealth[length=2.6mm]}] (a) -- node[gw, above=1pt]{blue $-1$} (b);
        \draw[gpath, -{Stealth[length=2.6mm]}] (b) -- node[gw, above=1pt]{red $\times2$} (c);
        \node[gv] (d) at (2,-1.3) {8};
        \node[gv] (e) at (4,-1.3) {7};
        \draw[gde] (a) -- node[gw, below left=-1pt]{$\times2$} (d);
        \draw[gde] (d) -- node[gw, below]{$-1$} (e);
        \draw[gde] (e) -- node[gw, right]{$-1$} (c);
        \node[note, anchor=west] at (4.5,-0.6) {orange: 2 presses\\gray: 3 presses};
    \end{scope}
\end{tikzpicture}
\figcap{Sample 1. Doubling first ($4\to8\to7\to6$) looks natural but costs 3 presses; stepping down first costs 2.}
\end{diagrambox}

\subsubsection*{2. The graph model}

\begin{modelbox}[Graph Model --- Two Buttons]
\renewcommand{\arraystretch}{1.35}
\begin{tabularx}{\linewidth}{@{}>{\bfseries}l X@{}}
Vertices & one vertex per number the display can show: $x \in \{1, 2, 3, \dots\}$; only the window $1 \le x \le \max(n, 2m)$ can matter (Lemma below) \\
Edges & directed: $x \to 2x$ (red) and $x \to x-1$ (blue, only if $x - 1 \ge 1$); generated on the fly, never stored \\
Weights & none --- every press costs 1 \\
Question & length of a shortest path from vertex $n$ to vertex $m$ \\
Size & $|V| \le 2\cdot10^4$, \ $|E| \le 4\cdot10^4$: BFS is instant \\
\end{tabularx}
\end{modelbox}

\begin{diagrambox}
\begin{tikzpicture}[x=1.25cm]
    \foreach \x in {1,...,9}
        \node[gv, minimum size=6.5mm] (n\x) at (\x,0) {\x};
    \node[font=\small, text=darkgray] at (10,0) {$\cdots$};
    \foreach \x in {2,...,9} {
        \pgfmathtruncatemacro{\y}{\x-1}
        \draw[gde, draw=primaryblue] (n\x) to[bend left=35] (n\y);
    }
    \draw[gde, draw=recursionred, line width=1.1pt] (n1) to[bend left=45] (n2);
    \draw[gde, draw=recursionred, line width=1.1pt] (n2) to[bend left=50] (n4);
    \draw[gde, draw=recursionred, line width=1.1pt] (n3) to[bend left=50] (n6);
    \draw[gde, draw=recursionred, line width=1.1pt] (n4) to[bend left=50] (n8);
    \node[font=\scriptsize\bfseries, text=recursionred] at (5.5,1.75) {red: $x \to 2x$};
    \node[font=\scriptsize\bfseries, text=primaryblue] at (5.5,-0.85) {blue: $x \to x-1$};
    \node[gvdead, minimum size=6.5mm] (z) at (0,0) {0};
    \draw[gcut] (0.5,-0.5) -- (0.5,0.5);
    \node[font=\scriptsize, text=warningred] at (0,-0.6) {breaks};
\end{tikzpicture}
\figcap{The implicit graph on the number line (only the first edges are drawn). Every vertex has at most two outgoing edges; the edge $1 \to 0$ is missing because the device would break.}
\end{diagrambox}

\subsubsection*{3. Key observations}

\begin{itemize}[leftmargin=*]
    \item \textbf{It is a shortest path with unit weights.} Each press is one edge, all presses cost the same, so BFS from $n$ gives the answer.
    \item \textbf{Forward greedy fails.} ``Double while below $m$, then step down'' gives $4\to8\to7\to6$ (3 presses) on sample 1; the optimum steps down \emph{first}. Choosing the right moment to subtract is exactly what BFS tries for us.
    \item \textbf{The graph is infinite} --- red can be pressed forever. Without a bound, BFS never stops. The lemma below gives the window $[1, \max(n, 2m)]$.
    \item \textbf{If $n > m$, red is useless:} the answer is $n - m$ (sample 2). This is the lemma applied at the very start.
\end{itemize}

\begin{insightbox}[{Lemma (the window) --- once the display shows $x \ge m$, red is never worth pressing}]
Take any way to go from $x \ge m$ to $m$ that uses $k \ge 1$ red presses, made at values $z_1, \dots, z_k \ge 1$, and $b$ blue presses. Each red press adds $z_i$ to the value, each blue press removes 1, so
\[
x + \sum_i z_i - b = m \quad\Longrightarrow\quad b = x - m + \sum_i z_i \;\ge\; x - m + k,
\]
and the total is $k + b \ge x - m + 2k > x - m$. But pressing blue $x - m$ times reaches $m$ directly (and never goes below $m \ge 1$). So in an optimal sequence, red is pressed only at values $x < m$, which produces at most $2(m-1)$; all other values come from $n$ by subtracting. \textbf{Every value on an optimal path is at most $\max(n, 2m-2)$.} \hfill$\blacksquare$
\end{insightbox}

\begin{diagrambox}
\begin{tikzpicture}[x=0.55cm]
    \fill[solutionbg] (1,-0.25) rectangle (20,0.25);
    \fill[recursionbg] (20,-0.25) rectangle (23.5,0.25);
    \draw[thick, -{Stealth}] (0,0) -- (23.7,0);
    \foreach \x/\l in {1/1, 6/m, 12/2m, 20/{\text{HI}=2\max(n,m)}} {
        \draw[thick] (\x,-0.3) -- (\x,0.3);
        \node[font=\scriptsize\bfseries, text=deepblue] at (\x,0.6) {$\l$};
    }
    \node[gvsrc, minimum size=5mm, font=\tiny\bfseries] (A) at (4,0) {$n$};
    \node[gvdst, minimum size=5mm, font=\tiny\bfseries] (B) at (6,0) {$m$};
    \draw[-{Stealth[length=2mm]}, recursionred, thick] (3,0.3) to[bend left=40] node[font=\tiny, above]{red at $x<m$} (6,0.3);
    \draw[-{Stealth[length=2mm]}, warningred, thick, dashed] (8,-0.3) to[bend right=30] node[font=\tiny, below]{red at $x \ge m$: never optimal} (16,-0.3);
    \node[font=\scriptsize, text=warningred] at (21.8,-0.6) {never needed};
\end{tikzpicture}
\figcap{The window when $n < m$: a red press made below $m$ lands below $2m$; a red press made at or above $m$ is never part of an optimal answer. The code uses $\text{HI} = 2\max(n, m) \ge \max(n, 2m)$.}
\end{diagrambox}

\subsubsection*{4. Approaches}

\begin{bruteforcebox}
\textbf{Idea: try every sequence of presses}, shortest first (iterative deepening). For $d = 0, 1, 2, \dots$ ask ``can $m$ be reached with at most $d$ presses?'' by trying red and blue at every step.

\begin{lstlisting}[style=cpp]
#include <bits/stdc++.h>
using namespace std;

// can the display go from x to m with at most k more presses?
bool canReach(long long x, long long m, int k) {
    if (x == m) return true;
    if (k == 0) return false;
    if (canReach(2 * x, m, k - 1)) return true;              // red
    return x - 1 >= 1 && canReach(x - 1, m, k - 1);          // blue (must stay positive)
}

int main() {
    long long n, m;
    cin >> n >> m;
    int d = 0;
    while (!canReach(n, m, d)) d++;                          // first depth that works
    cout << d << "\n";
}
\end{lstlisting}

\begin{diagrambox}
\begin{forest} dptree, for tree={font=\scriptsize\ttfamily}
[4, ans
  [8, edge label={node[midway, left, font=\tiny]{red}}
    [16] [7]]
  [3, edge label={node[midway, right, font=\tiny]{blue}}
    [6, bc] [2]]]
\end{forest}
\figcap{Depth 2 on sample 1: four sequences, one of them ends at 6 (green).}
\end{diagrambox}

\begin{complexitybox}
\textbf{Time:} $\Theta(2^{d})$ where $d$ is the answer. Correct on both samples, but $n = 10^4$, $m = 1$ needs $d = 9999$ --- about $2^{9999}$ sequences. The same number (say $7$) is re-explored along many different sequences. \textbf{Verdict:} \tle{}.
\end{complexitybox}
\end{bruteforcebox}

\begin{algorithmbox}
\textbf{Idea: BFS over the numbers $1..\text{HI}$.} The brute force re-explores the same number many times; BFS visits each number once, in order of the number of presses. By the Lemma, $\text{HI} = 2\max(n, m)$ is a safe window.

\begin{lstlisting}[style=cpp]
#include <bits/stdc++.h>
using namespace std;

int main() {
    int n, m;
    cin >> n >> m;
    const int HI = 2 * max(n, m);                  // every optimal path stays in [1, HI]
    vector<int> dist(HI + 1, -1);                  // dist[x] = fewest presses n -> x
    queue<int> q;
    dist[n] = 0;
    q.push(n);
    while (!q.empty()) {
        int x = q.front(); q.pop();
        if (x == m) break;                         // goal popped: distance is final
        int next[2] = {2 * x, x - 1};              // red, blue
        for (int y : next) {
            if (y < 1 || y > HI) continue;         // device breaks / outside the window
            if (dist[y] != -1) continue;           // already discovered
            dist[y] = dist[x] + 1;
            q.push(y);
        }
    }
    cout << dist[m] << "\n";
}
\end{lstlisting}

\begin{complexitybox}
\textbf{Time:} $\bigO{\text{HI}} = \bigO{n + m}$ --- at most $2\cdot10^4$ states, 2 edges each. \textbf{Space:} $\bigO{n + m}$. \textbf{Verdict:} \ac{}.
\end{complexitybox}
\end{algorithmbox}

\begin{invariantbox}[Why it is correct]
\textbf{Window:} by the Lemma some optimal sequence stays inside $[1, \text{HI}]$, so restricting BFS to the window removes no optimal path, and the window adds no invalid one (values $\ge 1$, and every edge is a real button press).
\textbf{Distances:} inside the window this is plain BFS with unit costs, so by the BFS invariant of Pattern~4 (\S\ref{ssec:ch1-p4-theory}) every \texttt{dist[y]} is set once, to the true fewest number of presses. $m$ is always reachable (blue from $n > m$; doubling past $m$ and stepping down from $n < m$), so \texttt{dist[m]} $\ne -1$.
\end{invariantbox}

\begin{alternativebox}[Approach 3: Reverse Greedy in $\bigO{\log m}$]
\textbf{Idea: run the film backwards.} Undoing red means \emph{halving} (only possible on an even number), undoing blue means \emph{adding 1}. Going backwards from $m$ to $n$, the choice becomes forced:
\begin{itemize}[leftmargin=*]
    \item $m$ odd $\Rightarrow$ it cannot be the result of a red press: add 1;
    \item $m$ even and $m > n$ $\Rightarrow$ halve (proof below);
    \item $m \le n$ $\Rightarrow$ the rest is $n - m$ blue presses (Lemma).
\end{itemize}
\begin{lstlisting}[style=cpp]
#include <bits/stdc++.h>
using namespace std;

int main() {
    int n, m;
    cin >> n >> m;
    int presses = 0;
    while (m > n) {                       // walk backwards from m
        if (m % 2 == 1) m += 1;           // undo a blue press
        else m /= 2;                      // undo a red press
        presses++;
    }
    cout << presses + (n - m) << "\n";    // the remaining part: blue only
}
\end{lstlisting}
\textbf{Why halving is right when $m > n$ is even.} Suppose an optimal backward sequence starts with $+1$. Plus-ones alone never get down to $n < m$, so some halving follows after $j \ge 1$ plus-ones, and $m + j$ must be even, so $j \ge 2$. Replace ``$j$ times $+1$, then halve'' ($j + 1$ steps) by ``halve, then $j/2$ times $+1$'' ($1 + j/2$ steps): both reach $(m + j)/2$, and the second is strictly shorter --- contradiction.

\begin{diagrambox}
\begin{tikzpicture}[node distance=1.9cm]
    \node[gvdst] (a) {10};
    \node[gv, right=of a] (b) {5};
    \node[gv, right=of b] (c) {6};
    \node[gvsrc, right=of c] (d) {3};
    \draw[back] (a) -- node[lbl, above]{even: $\div2$} (b);
    \draw[back] (b) -- node[lbl, above]{odd: $+1$} (c);
    \draw[back] (c) -- node[lbl, above]{even: $\div2$} (d);
    \node[note, below=3mm of b, xshift=1.9cm] {$n=3$, $m=10$: 3 backward steps, then $m = n$ (0 blue presses). Forward: $3 \xrightarrow{\times2} 6 \xrightarrow{-1} 5 \xrightarrow{\times2} 10$.};
\end{tikzpicture}
\end{diagrambox}

\begin{complexitybox}
\textbf{Time:} $\bigO{\log m}$ --- every two steps at least halve $m$. \textbf{Space:} $\bigO{1}$. \textbf{Verdict:} \ac{}; it also works for $m$ up to $10^{18}$, where BFS is impossible. The same task appears on LeetCode as ``Broken Calculator'' with values up to $10^9$ --- only the greedy survives there.
\end{complexitybox}
\end{alternativebox}

\begin{tricksbox}[Approach 4: Printing the Presses]
Contest problems often ask for the sequence itself. Store the parent when a number is discovered and walk back from $m$:
\begin{lstlisting}[style=cpp]
#include <bits/stdc++.h>
using namespace std;

int main() {
    int n, m;
    cin >> n >> m;
    const int HI = 2 * max(n, m);
    vector<int> dist(HI + 1, -1), par(HI + 1, -1);
    queue<int> q;
    dist[n] = 0;
    q.push(n);
    while (!q.empty()) {
        int x = q.front(); q.pop();
        int next[2] = {2 * x, x - 1};
        for (int y : next)
            if (y >= 1 && y <= HI && dist[y] == -1) {
                dist[y] = dist[x] + 1;
                par[y] = x;
                q.push(y);
            }
    }
    string presses;
    for (int v = m; v != n; v = par[v])
        presses += (par[v] * 2 == v ? 'R' : 'B');           // which button led to v?
    reverse(presses.begin(), presses.end());
    cout << dist[m] << " " << presses << "\n";               // "4 6" -> "2 BR"
}
\end{lstlisting}
The question ``was it red?'' is answered by $\texttt{par}[v] \cdot 2 = v$; red and blue can never lead from the same parent to the same $v$ (that would need $2x = x - 1$).
\end{tricksbox}

\subsubsection*{5. Trace}

\begin{tracebox}[{Sample 1: BFS from $n=4$ to $m=6$, window $[1, 12]$}]
\centering
\begin{tikzpicture}[x=1.15cm, y=0.75cm]
    \foreach \k/\px/\py/\title in {1/0/0/{pop 4}, 2/6.2/0/{pop 8}, 3/0/-4.4/{pop 3}, 4/6.2/-4.4/{pop 7, then pop 6}} {
        \begin{scope}[xshift=\px cm, yshift=\py cm]
            \node[font=\scriptsize\bfseries, text=deepblue, anchor=west] at (-0.4,1.9) {Step \k: \title};
        \end{scope}
    }
    % ---- step 1 ----
    \begin{scope}
        \node[gvcur] (a) at (0,0) {4};
        \node[gvseen] (b) at (1.3,1) {8};
        \node[gvseen] (c) at (1.3,-1) {3};
        \draw[gtree, -{Stealth[length=2.2mm]}] (a) -- (b);
        \draw[gtree, -{Stealth[length=2.2mm]}] (a) -- (c);
        \node[font=\scriptsize, anchor=west] at (2.3,0) {queue: \texttt{[8, 3]}};
    \end{scope}
    % ---- step 2 ----
    \begin{scope}[xshift=6.2cm]
        \node[gvdone] (a) at (0,0) {4};
        \node[gvcur] (b) at (1.3,1) {8};
        \node[gvseen] (c) at (1.3,-1) {3};
        \node[gvdead] (p) at (2.6,1.6) {16};
        \node[gvseen] (e) at (2.6,0.4) {7};
        \draw[gtree] (a) -- (b); \draw[gtree] (a) -- (c);
        \draw[gfaded, dashed] (b) -- (p);
        \draw[gtree, -{Stealth[length=2.2mm]}] (b) -- (e);
        \node[font=\scriptsize, anchor=west, align=left] at (3.2,-0.5) {16 $>$ HI: rejected\\queue: \texttt{[3, 7]}};
    \end{scope}
    % ---- step 3 ----
    \begin{scope}[yshift=-4.4cm]
        \node[gvdone] (a) at (0,0) {4};
        \node[gvdone] (b) at (1.3,1) {8};
        \node[gvcur] (c) at (1.3,-1) {3};
        \node[gvseen] (e) at (2.6,0.4) {7};
        \node[gvseen] (f) at (2.6,-0.6) {6};
        \node[gvseen] (g) at (2.6,-1.6) {2};
        \draw[gtree] (a) -- (b); \draw[gtree] (a) -- (c); \draw[gtree] (b) -- (e);
        \draw[gtree, -{Stealth[length=2.2mm]}] (c) -- (f);
        \draw[gtree, -{Stealth[length=2.2mm]}] (c) -- (g);
        \node[font=\scriptsize, anchor=west, align=left] at (3.2,-0.5) {6 discovered,\\dist 2\\queue: \texttt{[7, 6, 2]}};
    \end{scope}
    % ---- step 4 ----
    \begin{scope}[xshift=6.2cm, yshift=-4.4cm]
        \node[gvsrc] (a) at (0,0) {4};
        \node[gvdone] (b) at (1.3,1) {8};
        \node[gvdone] (c) at (1.3,-1) {3};
        \node[gvdone] (e) at (2.6,0.4) {7};
        \node[gvdst] (f) at (2.6,-0.6) {6};
        \node[gvseen] (g) at (2.6,-1.6) {2};
        \draw[gtree] (a) -- (b); \draw[gtree] (b) -- (e); \draw[gtree] (c) -- (g);
        \draw[gpath, -{Stealth[length=2.4mm]}] (a) -- (c);
        \draw[gpath, -{Stealth[length=2.4mm]}] (c) -- (f);
        \draw[gfaded, dashed, -{Stealth[length=2mm]}] (e) -- (f);
        \node[font=\scriptsize, anchor=west, align=left] at (3.2,-0.5) {$7\to6$: already seen\\6 popped: answer \textbf{2}\\path $4 \to 3 \to 6$};
    \end{scope}
\end{tikzpicture}

\begin{lstlisting}[style=trace]
pop  dist  red (2x)         blue (x-1)        queue after
 4    0    8  new, dist 1   3  new, dist 1    [8, 3]
 8    1    16 > HI, skip    7  new, dist 2    [3, 7]
 3    1    6  new, dist 2   2  new, dist 2    [7, 6, 2]
 7    2    14 > HI, skip    6  already seen   [6, 2]
 6    2    goal popped -> answer 2
\end{lstlisting}
\end{tracebox}

\begin{tracebox}[{Sample 2: $n = 10$, $m = 1$ --- BFS explores more than it needs}]
BFS also tries the red button: it discovers 20 numbers ($1..20$) before popping $1$ at distance 9. The layers are
\begin{center}\small
\begin{tabular}{@{}l*{10}{c}@{}}
\toprule
dist & 0 & 1 & 2 & 3 & 4 & 5 & 6 & 7 & 8 & 9 \\
numbers & 10 & 9, 20 & 8, 18, 19 & 7, 16, 17 & 6, 14, 15 & 5, 12, 13 & 4, 11 & 3 & 2 & 1 \\
\bottomrule
\end{tabular}
\end{center}
The reverse greedy needs no loop at all here ($m \le n$): the answer is $n - m = 9$ immediately.
\end{tracebox}

\subsubsection*{6. Test suite}

\begin{testbox}
\centering\small
\begin{tabular}{@{}cccl@{}}
\toprule
$n$ & $m$ & \textbf{Expected} & \textbf{Why} \\
\midrule
4 & 6 & 2 & sample 1; defeats ``double first'' \\
10 & 1 & 9 & sample 2; $n > m$: only blue \\
3 & 5 & 2 & $3\to6\to5$ goes \emph{above} $m$: kills the window $\text{HI} = m$ \\
7 & 13 & 2 & $7\to14\to13$, same trap \\
3 & 10 & 3 & reverse greedy uses both rules \\
1 & 10000 & 20 & long climb, many halvings backwards \\
10000 & 1 & 9999 & largest answer; brute force is hopeless \\
9999 & 10000 & 5000 & step down to 5000, then one red press \\
2 & 1 & 1 & smallest $n > m$ \\
1 & 2 & 1 & smallest $n < m$ \\
\bottomrule
\end{tabular}
\end{testbox}

All three solutions (iterative deepening on answers up to 14, BFS, reverse greedy) agree on every pair $1 \le n, m \le 200$, and BFS with the window agrees with BFS on $[1, 2\cdot10^5]$ and with the greedy on 3000 random pairs up to $10^4$.

\subsubsection*{7. Common mistakes}

\begin{warningbox}[Mistake 1 --- window too small]
\begin{lstlisting}[style=snippet]
const int HI = max(n, m);          // WRONG: n=3, m=5 needs 6; BFS prints -1
const int HI = 2 * max(n, m);      // RIGHT: Lemma
\end{lstlisting}
\end{warningbox}

\begin{warningbox}[Mistake 2 --- letting the display reach 0]
Without the test $y \ge 1$ the code reads \texttt{dist[0]} (and from there \texttt{dist[-1]}): undefined behaviour. Even if the array had room, 0 is a forbidden state --- the device breaks.
\end{warningbox}

\begin{warningbox}[Mistake 3 --- a greedy in the wrong direction]
Forward, neither ``double while below $m$'' nor ``subtract when it is odd'' is correct (sample 1). Backwards the choices are forced; forwards they are not. When a greedy fails in one direction, try the other.
\end{warningbox}

\begin{warningbox}[Mistake 4 --- forgetting the tail in the reverse greedy]
When the loop stops, $m$ may be \emph{below} $n$ (e.g.\ $n = 4$, $m = 6$: $6 \to 3$). The remaining $n - m$ blue presses must be added: $1 + (4 - 3) = 2$.
\end{warningbox}

\subsubsection*{8. Takeaways}

\begin{itemize}[leftmargin=*]
    \item ``Fewest operations from $n$ to $m$'' with unit-cost operations: numbers are vertices, operations are edges, BFS.
    \item An infinite state space needs a \emph{proven} window. Here: red pressed at $x \ge m$ never pays off, so $\text{HI} = \max(n, 2m)$.
    \item BFS is the safe general method; when the moves are invertible and become forced backwards, a reverse greedy can be exponentially faster.
    \item Store parents during BFS to print the actual sequence of moves.
\end{itemize}

\begin{variantsbox}
\variant{Multiply by 2 or append a 1}{\varin{https://codeforces.com/problemset/problem/727/A}}
  {red becomes $x \to 2x$, the second button becomes $x \to 10x + 1$; print any sequence from $a$ to $b$.}
  {same vertices (numbers up to $b$); both moves only increase, so the window is $[a, b]$.}
  {backwards, every number has at most one predecessor ($b/2$ if even, $(b-1)/10$ if it ends in 1): follow it (Exercise~E4.1).}
\variant{Three buttons}{\varnone}
  {a third button multiplies by 3.}
  {vertices $1..\max(n, 3m)$: the Lemma's argument works for any multiplication.}
  {BFS with three edges per vertex; the reverse greedy is no longer obviously optimal (Exercise~E4.2).}
\variant{Wrap-around display, many queries}{\varin{https://codeforces.com/problemset/problem/1661/B}}
  {operations $+1$ and $\times2$ modulo $32768$; fewest steps to reach 0 from each of many numbers.}
  {all $32768$ residues are vertices; the space is finite, no window needed.}
  {one BFS \emph{backwards} from 0 answers every query by a lookup (Exercise~E4.4).}
\variant{Buttons with different costs}{\varnone}
  {red costs $a$, blue costs $b$.}
  {the same graph, now with weights $a$ and $b$ on the edges.}
  {BFS no longer gives the cheapest sequence: Dijkstra's algorithm (Chapter 7).}
\end{variantsbox}
